\documentclass[11pt]{article}
\usepackage[a4paper,margin=25mm]{geometry}
\usepackage{amsmath,amssymb,amsthm,booktabs,xurl,hyperref}
\hypersetup{hidelinks}
\newtheorem{theorem}{Theorem}
\newtheorem{lemma}[theorem]{Lemma}
\newtheorem{proposition}[theorem]{Proposition}
\newcommand{\Q}{\mathbb Q}
\title{A quadratic ceiling-binary circuit for a fixed power}
\author{Z diophantine in Q project}
\date{23 November 2026}
\begin{document}
\maketitle
\begin{abstract}
A known nonzero base permits a power step of the form
$v'=v^2/b$ to be imposed by one quadratic equation.  Ceiling-binary
prefixes therefore compute an arbitrary fixed positive power with one
coordinate per binary position, including its subtraction step.  In the
combined Jones systems, the existing projective input square supplies
$b^3$ and $b^5$ in two coordinates.  Computing $b^{5^{60}}$ then takes
137 more coordinates, replacing the preceding 168-coordinate circuit.
The complete quartic counts fall from $1103,1283,837$ to
$1074,1254,808$.  We give full reconstruction maps, verify the nonzero
base before using division, and supply every resulting quadratic row.
This is an exact transformation of the preceding systems; their source
and arithmetic qualifications remain unchanged.
\end{abstract}

\section{The fixed-power circuit}

Let $M\ge1$ be an integer and put $L=\lceil\log_2 M\rceil$, with
$L=0$ when $M=1$.  Define integer exponents
\[
 e_j=\left\lceil\frac{M}{2^{L-j}}\right\rceil
 \quad(0\le j\le L),\qquad
 \delta_j=2e_{j-1}-e_j\quad(1\le j\le L).
\]
Then $e_0=1$, $e_L=M$, and every $\delta_j$ is either zero or one.
These statements are exact identities of integers, not estimates for a
numerical power evaluation.

\begin{theorem}\label{thm:power}
Let $K$ be a field and let $b\in K$ be nonzero.  Set $v_0=b$.
The following $L$ quadratic equations have a unique solution in
$v_1,\ldots,v_L$:
\begin{equation}\label{eq:circuit}
 \begin{cases}
 v_j-v_{j-1}^2=0,&\delta_j=0,\\
 bv_j-v_{j-1}^2=0,&\delta_j=1.
 \end{cases}
\end{equation}
That solution is $v_j=b^{e_j}$.  In particular its final coordinate is
$b^M$.  More generally, any already certified prefix $v_0=b^A$, with
$A=\lceil M/2^h\rceil$, starts the same construction with only $h$
new coordinates and equations.
\end{theorem}

\begin{proof}
For every real $t$, the difference
$2\lceil t\rceil-\lceil2t\rceil$ belongs to $\{0,1\}$.
Apply this with $t=M/2^{L-j+1}$.  Inductively, the first row of
\eqref{eq:circuit} gives $v_j=b^{2e_{j-1}}$, while the second gives
$v_j=b^{2e_{j-1}-1}$ because $b\ne0$.  These are exactly $b^{e_j}$.
Each row is linear in its new coordinate with nonzero coefficient
one or $b$, so the solution is unique.  The same argument applies
to the shorter sequence beginning at $A$.
\end{proof}

If nonvanishing of $b$ is not supplied by the surrounding system, the
extra equation $bt=1$ supplies it with one extra coordinate.  We do not
need that extra coordinate in the application below.  The theorem also
works for negative rational bases.  It asserts no representation of
$0^M$ by a chain that contains a division step.

\section{The base and prefix in the complete systems}

The retained outer variables include $b,B_5,g,t,m,q$, with coefficient
parameters $x,z,u,y$ in the positive-integer source.  Each positivity
wrapper has the form
\begin{equation}\label{eq:positive}
 2A-2=\sum_{j=1}^3 a_j(a_j+1).
\end{equation}
For rational witnesses, $a_j(a_j+1)\ge-1/4$, so every such wrapper
already implies $A\ge5/8>0$, before any integrality conclusion.
In particular the retained wrapper for $\theta=B_5-2z$ implies
$B_5>2z$.

In the general and specified interfaces write $c$ for the nonzero
common contact normalizer and $w_b$ for the stored input square.  The
literal input and unit equations imply
\[
 b^2=cw_b,\qquad u_c=c^{-1},
\]
where $u_c$ is an existing coordinate.  The unchanged two prefix rows
are
\begin{equation}\label{eq:prefix}
 u_c B_3=bw_b,\qquad u_c B_5=B_3w_b.
\end{equation}
Consequently $B_3=b^3$ and $B_5=b^5$, without division by $b$.
In the ordinary-$abc$ table the same equations have $c=u_c=1$.
Its input square is an unscaled coordinate $w_b=b^2$.
The coordinate called $B_3$ here has the retained machine name
\texttt{power\_001}; $B_5$ is called \texttt{B5}.

\begin{lemma}\label{lem:base}
For the retained positive integer parameters, every rational solution
of the unchanged input, unit, prefix, and positivity equations has
$b>0$.  This statement does not use any of the replaced power rows.
\end{lemma}
\begin{proof}
Equations \eqref{eq:prefix} give $B_5=b^5$ as just shown.  Since
$z\ge0$ and the $\theta$ wrapper gives $B_5>2z$, we have $b^5>0$,
and hence $b>0$.
\end{proof}

There is a second independent nonvanishing argument useful for other
power charts.  The unchanged wrappers give $g,t,m,\theta,\alpha>0$;
the defining rows give $e=y+m\theta>0$ and $\ell=u+t\theta>0$.
After undoing the existing invertible input-square rescaling, the
alpha row is
\[
 \alpha=(b-xy)q^2-e\ell g^2>0.
\]
It follows that $q\ne0$ and $b>xy$.  This argument likewise does not
assume a power relation between $q$ and $b$.  Thus the new construction
does not rely on a circular interpretation of its division steps.

\section{The 139-coordinate application}

Put
\[
 M=5^{60}=867361737988403547205962240695953369140625.
\]
Its binary length is 140, and
\[
 \left\lceil\frac{M}{2^{137}}\right\rceil=5.
\]
Keep the two old prefix coordinates $B_3,B_5$, and take
\[
 e_j=\left\lceil\frac{M}{2^{137-j}}\right\rceil,
 \quad 0\le j\le137.
\]
Beginning with $v_0=B_5$, Theorem~\ref{thm:power} adds 137 coordinates,
the last of which is the already named outer coordinate $q$.
The first exponents are $5,10,20,40,80,160$; the full recurrence,
including every choice of digit, is stored in the coefficient tables.
Exactly 78 of its 137 steps have $\delta_j=1$.

The old signed addition chain had 169 computed powers.  Its $b^2$
coordinate had already been removed using the integrality block's input
square, leaving 168 actual power coordinates.  The new circuit uses
$2+137=139$, including $B_3,B_5,q$.  In literal table terms it deletes
165 private intermediate coordinates and adds 136, for a saving of 29.
It deletes 166 rows and adds 137, also saving 29 rows.  The two unchanged
prefix rows are counted only once.

\begin{theorem}\label{thm:application}
For each admissible fixed coefficient choice, the complete preceding
systems and the new systems have identical rational solution projections
to all coordinates outside the power circuit.  The new counts are
\begin{center}
\begin{tabular}{lrrrr}
\toprule
Interface & Old variables & New variables & Old rows & New rows\\
\midrule
General &1103&1074&953&924\\
Specified &1283&1254&1109&1080\\
Ordinary rational $abc$ &837&808&725&696\\
\bottomrule
\end{tabular}
\end{center}
The sums of squares of the new residuals have total degree exactly four.
The variable counts exclude the integer instance encoded in coefficients.
\end{theorem}

\begin{proof}
By Lemma~\ref{lem:base}, both systems have $b\ne0$ before interpreting
either circuit.  The old chain is a sequence of multiplication steps
and divisions by already computed positive powers of $b$.  Its complete
exponent certificate gives $q=b^M$.  Theorem~\ref{thm:power} gives the
same conclusion for the new chain.  Starting with a solution of either
system, compute each removed power as its specified integer power of
the same $b$.  This recovers all rows of the other system.  Every other
coordinate and every other residual is preserved literally, up to the
renumbering of coordinates.

The full occurrence audit proves that no deleted intermediate occurs
outside a deleted power row.  This makes the substitution valid for the
complete systems, including the scalar, contact, positivity, and Pell
blocks.  The two coordinate maps are inverse on their solution sets,
since both circuits have uniquely determined intermediates at fixed $b$.

Each new row has degree at most two and integer coefficients.  A sum
of rational squares is zero exactly when each summand is zero.  At least
one unchanged outer row has a nonzero quadratic homogeneous part, so the
sum of squares has degree exactly four: the sum of squares of the leading
quadratic parts is a nonzero polynomial over the ordered field $\Q$.
The direct census above gives the displayed counts.
\end{proof}

\section{A restricted optimality statement}

The numerical count is not a general lower bound for polynomial
encodings.  There is, however, a simple lower bound for the exact
positive-power chart used here.

\begin{proposition}
Suppose the initially available monomial coordinates have base
exponents at most two and all base exponents are nonnegative.
Each subsequent coordinate is determined linearly by a quadratic
binomial equation whose other coordinates are already available.
Then computing the power $b^M$, for $M=5^{60}$, requires at least
139 subsequent coordinates.  The present 139-coordinate circuit
attains this bound within that chart.
\end{proposition}

\begin{proof}
Assign exponent zero to the independent common normalizers.
If $E$ is the largest available exponent, the new exponent has the
form $a+b-d$ with $0\le a,b,d\le E$, allowing an omitted factor
to have exponent zero.  It is at most $2E$.  Thus after $t$ new
coordinates the largest exponent is at most $2^{t+1}$.  Since
$2^{139}<5^{60}<2^{140}$, one needs $t\ge139$.
The projective input square has base exponent two, and the unchanged
prefix \eqref{eq:prefix} and every new ceiling step satisfy the stated
rules with nonnegative exponents.  There are exactly 139 new power
coordinates relative to that initial pool.
\end{proof}

The proposition does not cover negative-power coordinates, nontriangular
systems, other algebraic variables, Pell certificates, changes to the
universal source, or the use of source equations to avoid computing this
power.  Those possibilities remain open.

\section{Reproducibility and qualifications}

The complete tables and certificates are in
\nolinkurl{undecidability/coefficients/jones1980_ceiling_power/}.
The builder is
\nolinkurl{undecidability/code/build_ceiling_power.py}; it also accepts
an explicit predecessor and output path for composition with independent
changes.  The independent checker
\nolinkurl{undecidability/verify_ceiling_power.py} reconstructs the exponent
recurrence from the actual new equations, checks the old and new
coordinate occurrences, compares every unchanged residual, verifies
integer parameter-polynomial coefficients and degree, and reproduces
the canonical files byte for byte.  Its small rational examples supplement
the exact exponent certificate; they do not replace it.

The first two predecessors are the detector-trace and Pell systems
\cite{detector,pell}; the third is the shared-reference system after
the Pell reduction.  Their precise emended source universality premise
$U_{\rm em}$ is retained in these historical tables.  In particular this
transformation does not correct the printed $n^2$ binomial parameter,
nor assert universality of a corrected source.  A separate corrected
source may be composed with this transformation after its proof is
supplied.  Every retained integer-predicate arithmetic premise remains
as stated in the predecessor papers.  The ordinary-$abc$ qualification
applies to the third row only.  None of these algebraic transformations
claims an unconditional solution of Hilbert's tenth problem over $\Q$,
or a quartic bound near 263.

\begin{thebibliography}{9}
\bibitem{detector} Z diophantine in Q project,
\emph{A rational trace in place of each private detector product},
\nolinkurl{undecidability/papers/detector_trace_compression.tex}.
\bibitem{pell} Z diophantine in Q project,
\emph{Pell index optimization},
\nolinkurl{undecidability/papers/pell_index_optimization.tex}.
\end{thebibliography}
\end{document}
